% Partial manuscript section. Not a complete manuscript or novelty assertion.
% Requires amsmath, amssymb, amsthm; theorem/proposition environments supplied by host.
\section{Grounded budget certificates for cyclic extraction}
Let $G$ be an explicitly represented equality graph with root set $R$ and
nonnegative integer node costs. An extraction selects one representative for
each used class and must induce an acyclic computation. Each shared result is
charged once. All equivalences are upstream assumptions.
Let $C$ contain the roots and every potentially reachable class referenced by
at least two distinct parent classes; write $s=|C|$. A private class has a
unique boundary owner and the private support is acyclic. For $S\subseteq C$,
let $c_v(S)$ be the minimum cost of a private local expansion constructing $v$
using only $S$ as available boundary results. Unavailable expansions cost
$+\infty$. The local functions are computed in polynomial time by min-plus
evaluation of the private support.

For $b\in\mathbb Z_{\ge0}^C$, set
\[
 S_0(b)=\varnothing,\qquad
 S_{t+1}(b)=S_t(b)\cup\{v\in C:c_v(S_t(b))\le b_v\}.
\]
Since $c_v$ is nonincreasing in its argument, this reaches a fixed point
$S_*(b)$ after at most $s$ changing rounds.

\begin{theorem}[Grounded budget certificate]
For every integer $K\ge0$, the following are equivalent:
\begin{enumerate}
\item $G$ has an acyclic extraction containing $R$ with cost at most $K$.
\item There is $b\in\mathbb Z_{\ge0}^C$ with $\sum_v b_v\le K$ and
$R\subseteq S_*(b)$.
\end{enumerate}
\end{theorem}
\begin{proof}
For the second-to-first implication, retain a minimizing local witness when
$v$ first enters the closure. Its boundary dependencies were constructed in
earlier rounds, and private support is acyclic. The concatenated witnesses
therefore give an acyclic program. Their private pieces have distinct owners,
so their costs add without double counting. Each is charged at most $b_v$.
Pruning components unused by the roots cannot increase nonnegative cost.

Conversely, partition an acyclic extraction into its boundary-owned local
pieces. Assign to each used $v$ its actual local cost and assign zero to
unused boundary classes. The resulting vector sums to the extraction cost.
Induction on any dependency-first ordering of the used boundary classes
shows that each class enters the closure: once its predecessors are present,
the definition of $c_v$ admits the extraction's piece at cost at most $b_v$.
This includes zero-cost pieces; a zero cap is not an assumption of availability.
\end{proof}

\begin{corollary}[Circuit-model upper bound; comparison qualified]
The threshold problem admits a bounded-error quantum algorithm using
\[
 O\!\left(\sqrt{\binom{K+s}{s}}\,
       \operatorname{poly}(N,L,\log(K+2))\right)
\]
time and polynomial workspace, where $N$ is explicit graph size and $L$
includes the binary input-cost lengths. No quantum random-access table of
exponential size is required.
\end{corollary}
\begin{proof}
There are $M=\binom{K+s}{s}$ nonnegative cap vectors with sum at most $K$.
A rank in $[0,M)$ can be decoded with binomial prefix counts and binary
search; the computation uses polynomial time and space in $s$ and
$\log(K+2)$, not a loop of length $K$. Prepare a uniform superposition of
$q=\lceil\log_2M\rceil$ bits, rejecting padded indices. Evaluate the grounded
certificate reversibly, mark acceptance, and uncompute. Saturate local costs
at $K+1$, so all arithmetic has polynomial bit width. Standard search for
an unknown number of marked elements uses $O(\sqrt M)$ evaluations at
constant bounded error. Decode the measured rank, recover the local
witnesses, prune, and verify the classical output. When $M=1$, evaluate it
directly. The reversible construction follows ordinary computation-history
storage and uncomputation; its workspace is polynomial in the explicit
instance. The implementation accompanying this section compiles the binary
cap marker, but not the complete compact-index decoder and search circuit.
\end{proof}

This corollary improves allocation enumeration, not necessarily the best
classical algorithm. Classical subset dynamic programming has a
$2^s\operatorname{poly}(N,L)$ bound; representative enumeration, treewidth
methods, memoization and heuristics may be better. Numeric rescaling of
costs changes the cap domain without changing the extraction problem.
Relevant cost granularity and domain redundancies must therefore be included
in any significance or application claim. The fixed-point, min-plus and
quantum-search tools are prior techniques; priority for this structural
combination remains to be established.
